\subsection{拓扑群的连通分支}

设 V 是 e 在 G 中的一个对称邻域。由 V 生成的子群记作 \(V^{\infty}\)，它由 V 中元素的有限序列的所有乘积 \(\prod_{i=1}^{n} x_{i}\) 组成。\(V^{\infty}\) 是开的，因为它以 e 为一个内点，因而由第 1 目命题 4 是闭的。由此可得：

命题 6 连通拓扑群由单位元的每个邻域生成。

正如我们将在第四章（§ 2 第 5 目）看到的，这个命题的逆命题一般不成立。若拓扑群 G 由单位元的每个邻域生成，我们至多能说 G 除 G 外不含任何开子群。

* 作为具有异于 G 的开子群的非连通群 G 的例子，可以举出非零实数的乘法群 \(\mathbf{R}^{*}\)，其中子群 \(\mathbf{R}_{+}^{*}\)（由实数 \textgreater0 组成的）既开又闭（见第四章 § 3 第 2 目）。*

命题 7 在拓扑群 G 中，单位元 e 的连通分支 K 是一个闭正规子群。任一点 \(x \in G\) 的连通分支是陪集 \(x.K = K.x\)。

若 \(a \in \mathbf{K}\)，则 \(a^{-1}\mathbf{K}\) 连通且包含 \(e\)；于是 \(\mathbf{K}^{-1}\mathbf{K} \subset \mathbf{K}\)，这表明 \(\mathbf{K}\) 是 \(\mathbf{G}\) 的一个子群。这个子群在 \(\mathbf{G}\) 的所有自同构下不变，特别在所有内自同构下不变，因此 \(\mathbf{K}\) 在 \(\mathbf{G}\) 中是正规的；又 \(\mathbf{K}\) 是闭的（第一章 § 11 第 5 目命题 9）。最后，左平移 \(y \to xy\) 是 \(\mathbf{G}\) 的一个把 \(e\) 映为 \(x\) 的同胚，因而 \(x\) 的连通分支是 \(x\cdot \mathbf{K}\)。

G 的单位元 e 的连通分支称为 G 的单位元连通分支。

